\section{What still obstructs a general quasi-polynomial bound}\label{sec:barriers}

\subsection{Finite matching types do not bound multiplicity}

The macro state space separates geometry from algebraic copies.
Even a very small collection of boundary matchings can occur with
exponentially large multiplicities. Streaming changes how those copies
are visited; it does not eliminate the copies. The one-run recurrence
does eliminate a repeated algebraic region, but only after a fixed
finite context has been solved.

The three-strand obstruction in Kelom\"aki--Sch\"utz~\cite{ks26}
is particularly decisive against a universal width-only argument.
Some constant-strand families have exponentially growing total Betti
number, so any algorithm explicitly retaining a basis must expand
exponentially. Their specialized polynomial computation shows why this
is a representation obstruction rather than an impossibility theorem
for compact output.
The same lesson applies to scalar transfer matrices: their coefficients
can encode very large counts, but a transfer formula for the Euler
characteristic does not automatically determine homology ranks.

Likewise, the restriction \(t=O(\log n)\) in
\eqref{eq:prior-bound} is substantive. It cannot be removed by calling
the supplied run count a width parameter. A general diagram may have
\(t=\Theta(n)\), and the existence of some small topological
decomposition does not imply a polynomial procedure to find it with
short, verifiable moves.

\subsection{Why several long variables need more than interpolation}

A plausible next direction is to allow several large exponents at
once. Even the simplest example rules out one overly strong formulation.

\begin{proposition}[No single eventual polynomial in all twist lengths]
\label{prop:chambers}
There is no polynomial \(p(m,\ell)\) agreeing, for all sufficiently
large \(m,\ell\) with \(m-\ell\) odd, with the total reduced
characteristic-two Khovanov rank of the closed braid
\(\sigma_1^m\sigma_1^{-\ell}\).
\end{proposition}

\begin{proof}
The braid cancels to \(\sigma_1^{m-\ell}\), a two-strand torus knot
when \(m-\ell\) is odd. Its reduced rank is \(|m-\ell|\).
This last fact also follows directly from the one-run normal form:
there is one class in degree zero, none in degree of magnitude one,
and one in each subsequent degree through magnitude \(|m-\ell|\).

If \(p\) existed, fix any sufficiently large integer \(\ell\).
For infinitely many admissible \(m>\ell\),
\(p(m,\ell)=m-\ell\). A one-variable polynomial with infinitely many
zeros is zero, so \(p(x,\ell)=x-\ell\) identically in \(x\).
This holds for infinitely many \(\ell\), hence
\(p(x,y)=x-y\) as a two-variable polynomial.
It contradicts the required value \(\ell-m\) for admissible
\(\ell>m\).
\end{proof}

The cancellation in this example is cheap and should be performed first
in a practical recognizer. Its role is conceptual: comparable exponents
can produce different asymptotic chambers separated by equalities of
lengths. A multivariable theorem might use piecewise polynomial,
quasi-polynomial, or rational generating descriptions with explicit
chamber conditions. A few numerical values cannot select the correct
description or prove its validity throughout a large orthant.

\subsection{A sufficient representation target}

The useful target is a compact algebraic model closed under the operations
the algorithm actually performs. For a tangle fragment it must include
boundary data, grading shifts, multiplicities, differentials, and the
maps used in subsequent gluing. A list of dimensions is insufficient.

\Needspace{14\baselineskip}
\begin{proposition}[Conditional complexity contract]\label{prop:contract}
Suppose every \(n\)-crossing input admits a constructible decomposition
with \(\poly(n)\) composition steps. Assume:
\begin{enumerate}
\item each intermediate fragment has a faithful algebraic description
of at most \(S(n)\) bits, including all multiplicities and map data;
\item each composition, reduction, and final rank-one decision can be
performed in \(\poly(n,S(n))\) bit operations while preserving a
description of that size;
\item the construction and the correctness of every representation
change have the same total complexity bound.
\end{enumerate}
If \(S(n)=2^{O((\log n)^2)}\), this gives a complete
\(n^{O(\log n)}\)-time recognition algorithm.
\end{proposition}

\begin{proof}
There are polynomially many steps, each of polynomial degree independent
of \(n\) in \(n\) and \(S(n)\). Their sum, including the construction and
final decision, remains \(2^{O((\log n)^2)}\).
Faithfulness makes the final rank-one decision equivalent to
\eqref{eq:detector}.
\end{proof}

This proposition is a contract for future theorems, not a proof that such
descriptions exist. In particular the first assumption forbids hiding
large basis multiplicity inside a supposedly small number of geometric
types, and the second forbids an implicit exponential expansion during
composition. The third distinguishes finding a useful decomposition
from merely asserting its existence.

The present recurrence satisfies this kind of contract for one finite
context followed by a homogeneous repeated chain region. It gives a
concrete example of how much information can be compressed while still
recovering exact ranks. Extending that closure property to interacting
regions is the central algebraic problem.

\section{Further research questions and proposed work programme}
\label{sec:research}

The following topics are ordered by their connection to the actual
remaining bottlenecks. Each has a concrete mathematical or experimental
completion criterion. They are proposals, not claimed consequences of
the current implementation.

\subsection{Reduce the effective context amplitude}

\begin{question}\label{q:amplitude}
Can the raw crossing span \(W\) be replaced by the degree span of a
smaller context model that retains both dot actions and the distinguished
saddle?
\end{question}

The proof of \cref{thm:recurrence} uses a common support interval for the
\(I\) and \(E\) context sectors and compatibility with their maps.
Suppose an explicitly constructed equivalence of this two-sector
context has degree-preserving forward and backward maps, with contracting
homotopies, all compatible with both exterior dot actions and the
distinguished saddle. If both resulting sector complexes are supported
in \([a',c']\), the same proof gives threshold \(c'-a'+2\).
Compatibility is required at the level of maps and homotopies.
An ordinary homotopy equivalence of the underlying complexes alone
does not suffice; transporting operators through an arbitrary
contraction can introduce additional terms.

A first implementation should perform local contractions in the open
tangle category and transport the saddle, dot actions, and contracting
homotopies. It should output an equivalence witness and measure both
the reduced amplitude and the resulting matrix size.
A useful negative outcome would be a family where amplitude is small
but the equivariant context dimension still grows exponentially.
That would isolate the next barrier more accurately than crossing
count alone.

\subsection{Two interacting variable regions before arbitrary many}

\begin{question}\label{q:two}
For two distinguished runs, can their simultaneous long region be
described by finitely many exact face and chamber models whose evaluation
is polynomial in the two exponent bit lengths?
\end{question}

The one-run interior is controlled by \(D_E+L\). With two selected
turnbacks, there are two commuting local dot sums and several
\(I/E\) sectors. The total homological degree combines both local
coordinates, so the number of copies along a diagonal changes near
the boundary of the exponent rectangle. A faithful model must retain
those incidences, not simply add two independently computed slopes.

Begin with adjacent and nonadjacent generators, both relative signs,
all original component parities, and the cancellation family in
\cref{prop:chambers}. Require an exact symbolic formula whose chamber
conditions are proved, then compare complete small profiles with the
crossing cube. Only after the two-variable case is understood should
one study an \(s\)-variable face system with up to \(2^s\) sectors.
Even that sector count would not bound internal module complexity.

\subsection{Restore quantum information}

\begin{question}\label{q:quantum}
What is the exact bigraded refinement of the one-reference recurrence,
including the shifts of both boundary regions and the repeated interval?
\end{question}

The ungraded deformation \(D_E+L\) combines maps with different original
bidegrees. The slope identity gives total dimension and does not place
the smoothing's original bigraded groups in the long complex.
A correct refinement must carry the quantum shifts in the integer-twist
normal form before flattening the central groups.
One possible output is a rational generating expression in homological
and quantum variables, with finite boundary corrections. Its validity
and output size have to be proved.
The completion criterion is agreement with an independently bigraded
cube on both signs and mixed contexts, not just recovery of the Jones
Euler polynomial.

\subsection{Odd characteristics and integral torsion}

\begin{question}\label{q:integral}
Can a signed periodic central model replace \(F\) over fields of odd
characteristic and over the integers?
\end{question}

The current square-zero calculation and weight cancellation use
characteristic two. Over the integers, twist differentials alternate
dot sums and differences; signs also interact with tensor degrees.
A period-two central model is a natural object to examine, but using
the present \(D_E+L\) unchanged is invalid.
An integral algorithm additionally needs Smith-form or module data,
with the bit lengths of torsion invariants included in its complexity.
The first result should be a fully signed local and contextual theorem,
followed by a recurrence for isomorphism types or a precisely weaker
rank statement.

\subsection{Decision certificates smaller than full homology}

\begin{question}\label{q:certificate}
Can a rank-greater-than-one witness be represented and verified compactly
on inputs that evade the structural and modular filters?
\end{question}

There is a concrete algebraic witness to seek. In a chain complex with
total differential \(d\), find two cycles \(z_1,z_2\) and two linear
functionals \(\ell_1,\ell_2\) satisfying
\[
dz_j=0,\qquad \ell_i d=0,\qquad \ell_i(z_j)=\delta_{ij}.
\]
Then the two homology classes are independent: any boundary relation
between them is annihilated by both functionals. This elementary
criterion gives a verifiable target.
An explicit witness can be exponentially long, so the research problem
is a representation that supports these products without expansion.
The streamed lower bound already proves a decision but does not yet
produce such a portable compressed cycle--cocycle witness.

Unknot certificates need a different positive mechanism, such as a
small contraction to the one-generator reduced complex or a
topological simplification certificate. A method that quickly produces
two surviving classes on many knots still needs a complete, costed
endpoint on hard unknots.

\subsection{Use tangle invariants with explicit algebraic size accounting}

\begin{question}\label{q:curves}
Can immersed-curve or related four-ended tangle descriptions support
exact repeated gluing with compact multiplicities and controlled local
systems?
\end{question}

The framework of Kotelskiy--Watson--Zibrowius is a serious source of
structure~\cite{kwz}. The relevant complexity parameters would include
curve segments, intersections under twisting, local-system ranks and
entries, and the cost of cancelling the resulting Floer-type complex.
A picture with few curves may still encode many intersections or a
large local system. The first milestone should translate one delivered
long-run example into that representation and recover its threshold,
slope, and boundary corrections with explicit cost accounting.

\subsection{Recover speed while preserving the streamed memory bound}

\begin{question}\label{q:streamspeed}
Which bounded low-level changes recover the streamed backend's observed
time loss without restoring unbounded tables?
\end{question}

Candidates include a compiled exact bitset kernel, small-label tables
admitted only under an explicit slot cap, reused neighboring state keys,
and a more efficient geometry construction.
The benchmark currently shows a real time cost, so a default switch
would be premature.
Any follow-up should report complete output equality, paired A/B and
A/A timings, total versus live state counts, map-compilation counts,
and separate traced-allocation measurements on the same inputs.

A related question is whether the sweep should start at the other
homological end. Reversing the sweep requires the correct transposed
local maps and grading conversion. A preflight estimate may choose an
end, but an estimate of where homology is likely cannot substitute
for the exact incident ranks required by an early decision.

\subsection{Benchmark the unresolved part of the pipeline}

\begin{question}\label{q:corpus}
Which checked inputs actually reach the expensive endpoint after all
current inexpensive certificates, and what algebraic feature predicts
their cost?
\end{question}

Record the exit stage of every case, rather than timing all backends
on examples already decided at the front. Include actual hard unknot
fixtures, determinant-one nontrivial knots, mixed-sign braids of at
least four strands, and examples with large intermediate but small
final homology. Retain source provenance and the distinction between
PD-only and braid-source inputs.
The correct metrics include crossings, runs, degree span, maximum
layer dimension, cancellation fill, map recompilation, and the
compressed context size from \cref{q:amplitude}.
Artificially labeling an easy torus or three-strand example as hard
would make this research programme less informative.

\subsection{Cost the hierarchy route operation by operation}

\begin{question}\label{q:hierarchy}
Can the remaining hierarchy steps be given explicit input/output bit
contracts and composed into a complete quasi-polynomial runtime proof?
\end{question}

The required ledger includes hierarchy length, boundary-pattern size,
surface-coordinate bit length, component tracking after cutting,
product-region recognition, and the search cost of each next surface.
Polynomial certificate verification must be kept separate from finding
the certificate. Compact normal coordinates must remain compact through
every update.
Lackenby's current results provide substantial ingredients
\cite{lackenby26}; completing this ledger would be an independent route,
not a corollary of the present Khovanov recurrence.

\subsection{Formalize the finite algebraic core first}

\begin{question}\label{q:formal}
Which parts admit a small proof-assistant development with a stable
boundary to the topological literature?
\end{question}

The square-zero calculation, signed interval bookkeeping, rank identity,
positive and negative profile formulas, and cycle--cocycle witness
criterion are finite algebra and arithmetic. They are natural first
targets. A formalization should expose the local tangle-normal-form
assumption and then separately formalize or import its topological
proof. A machine-checked recurrence conditional on that assumption is
valuable, but is not a formalization of the complete unknot detector.
The current package contains conventional proofs and executable
cross-checks, not a Lean proof.

\subsection{A proposed sequence of research milestones}

\begin{enumerate}
\item Integrate the additive code and preserve the default dispatcher.
Reproduce the 254-method suite and the independent cube checks.
\item Build dot- and saddle-compatible context contractions and measure
effective amplitude on filter-undecided examples.
\item Prove a complete two-variable theorem, including its chambers,
component parities, and compact output.
\item Construct and independently verify compact nontriviality witnesses,
with an explicit positive endpoint for unknot cases.
\item Test a representation closed under repeated tangle gluing against
the size and operation contract in \cref{prop:contract}.
\item Pursue the hierarchy cost ledger in parallel, so the programme
does not depend on one algebraic representation succeeding universally.
\end{enumerate}

The immediate work is concrete and testable. A general quasi-polynomial
claim should be made only when both a compact representation and the
cost of every operation on it have been proved, including construction
and the final decision.
